\documentclass[11pt]{article}
\usepackage{ghb}

\title{Counting triangles by degeneracy ordering}
\date{2026-08-16}
\summary{Why orienting a graph by its degeneracy ordering turns triangle
counting into a near-linear problem, and what the bound actually depends on.}

\begin{document}

\begin{abstract}
Counting triangles naively costs $\bigO{n^3}$. Orienting the edges along a
degeneracy ordering brings it down to $\bigO{m \cdot d}$, where $d$ is the
degeneracy --- a quantity that is small for most graphs people actually care
about. This post works through why.
\end{abstract}

\section{The problem}

Given an undirected graph $G = (V, E)$ with $\abs{V} = n$ and $\abs{E} = m$, we
want $T(G)$, the number of triangles. Checking every triple is $\bigO{n^3}$,
which is hopeless past a few thousand vertices.

The useful observation is that a triangle has three vertices, and if we impose
\emph{any} total order on $V$, exactly one of them is smallest. So if every
triangle is counted from its lowest-ordered vertex, each is counted once.

\section{Orienting the graph}

\begin{definition}[Degeneracy]
The \emph{degeneracy} $d$ of $G$ is the smallest $k$ such that every subgraph
of $G$ has a vertex of degree at most $k$.
\end{definition}

A degeneracy ordering is produced greedily: repeatedly remove a
minimum-degree vertex, and record the order of removal. Orienting each edge
from earlier to later in that ordering gives a directed acyclic graph in which
every out-degree is at most $d$.

\begin{tikzpicture}[node distance=15mm]
  \node[vertex] (a) {$1$};
  \node[vertex, right=of a] (b) {$2$};
  \node[vertex, right=of b] (c) {$3$};
  \node[vertex, below=11mm of b] (d) {$4$};

  \draw[arc] (a) -- (b);
  \draw[arc] (b) -- (c);
  \draw[arc] (a) -- (d);
  \draw[arc] (b) -- (d);
  \draw[arc] (c) -- (d);
  \draw[arc] (a) to[bend left=32] (c);
\end{tikzpicture}

Every edge points from the lower number to the higher one, so no cycle can
close and the orientation is acyclic.

\section{The bound}

\begin{theorem}[Chiba--Nishizeki]
\label{thm:main}
Triangles in $G$ can be counted in time $\bigO{m \cdot d}$, where $d$ is the
degeneracy of $G$.
\end{theorem}

\begin{proof}
Orient $G$ by a degeneracy ordering, so every out-degree is at most $d$. For
each vertex $u$, examine every pair of out-neighbours $v, w$ and test whether
the edge $vw$ is present. Each triangle is found exactly once, from its
lowest-ordered vertex.

The work at $u$ is $\bigO{\deg^{+}(u)^2}$, so the total is
\begin{equation}
  \sum_{u \in V} \deg^{+}(u)^2
    \;\le\; \max_{u} \deg^{+}(u) \cdot \sum_{u \in V} \deg^{+}(u)
    \;\le\; d \cdot m,
  \label{eq:sum}
\end{equation}
using $\sum_u \deg^{+}(u) = m$ and $\deg^{+}(u) \le d$. Adjacency tests are
$\bigO{1}$ with a hash set.
\end{proof}

\begin{remark}
The bound is driven by degeneracy, not by $n$. Planar graphs have $d \le 5$,
and most sparse real-world networks sit in the tens, so
Theorem~\ref{thm:main} is close to linear in practice even when the maximum
degree is enormous.
\end{remark}

\section{What it costs in practice}

\begin{center}
\begin{tabular}{lrrr}
\hline
Graph & $n$ & $m$ & $d$ \\
\hline
Planar (any)      & ---       & $\le 3n-6$  & $\le 5$ \\
Road network      & 2.0M      & 2.8M        & 3 \\
Social graph      & 4.0M      & 88M         & 115 \\
\hline
\end{tabular}
\end{center}

The middle column is what a naive bound would key off; the last column is what
actually governs the running time.

\section{Related work}

The $\bigO{m \cdot d}$ result is due to Chiba and
Nishizeki~\cite{chiba1985arboricity}, stated there in terms of arboricity,
which is within a constant factor of degeneracy. Matrix-multiplication methods
do better asymptotically for dense graphs~\cite{alon1997finding}, and
Schank and Wagner~\cite{schank2005finding} give a careful experimental
comparison.

\bibliographystyle{plain}
\bibliography{refs}

\end{document}
